% ============================================================
% Appendix Z.12 — FRFP Case Study (Platforms / Everyday Use):
% Recommenders, Machine Translation, and ChatGPT
% ============================================================

\section{FRFP Case Study (Platforms / Everyday Use):
Recommenders, Machine Translation, and ChatGPT}
\label{app:frfp-platforms}

This appendix treats three families of widely deployed AI systems as a single
FRFP-style case study:

\begin{enumerate}
\item large-scale recommender systems (e.g.\ YouTube, Netflix, Amazon);
\item neural machine translation (NMT), exemplified by Google Translate;
\item large language models (LLMs) in general-purpose tools like ChatGPT.
\end{enumerate}

These systems are striking AI successes:

\begin{itemize}
\item Recommenders and ranking models drive much of the engagement and
revenue of modern platforms, and increasingly mediate what people
watch, read, and buy.
\item NMT (e.g.\ Google Neural Machine Translation, GNMT) has moved web
translation from brittle phrase-based systems to fluent, sentence-
level neural models across 100+ languages.
\item ChatGPT and similar LLM tools now have hundreds of millions of weekly
users and process billions of prompts per day, becoming a general
“language layer” for everyday tasks.
\end{itemize}

Unlike AlphaFold or DR screening AI, these systems live directly in the
everyday cognitive environment of billions of people, which makes FRFP’s
explicit/tacit and population-level lenses especially important.

\subsection{Explicit and Tacit Spaces for Everyday AI Platforms}

\subsubsection*{Explicit space (E_{\mathrm{platform}})}

Across these systems, explicit artefacts include:

\begin{itemize}
\item user interaction logs:
clicks, watch time, purchases, queries, prompts, translations;
\item content and item representations:
videos, products, web pages, sentences, embeddings, token sequences;
\item model parameters:
recommender networks, translation models, LLM weights;
\item outputs:
ranked lists of items, translated text, generated responses;
\item metrics:
CTR, dwell time, retention, BLEU/COMET scores, engagement,
satisfaction proxies.
\end{itemize}

These are all explicit structures in (E_{\mathrm{platform}}), manipulated
by training and serving pipelines in (N \ltimes E).

\subsubsection*{Tacit space (T_{\mathrm{platform}})}

Tacit space contains:

\begin{itemize}
\item users’ beliefs, preferences, habits, and values;
\item creators’ and merchants’ strategies in responding to platform
incentives;
\item engineers’ and product managers’ intuitions about trade-offs
(e.g.\ diversity vs.\ engagement, safety vs.\ expressivity);
\item regulators’ and the public’s evolving norms about fairness, safety,
and acceptable influence;
\item workers’ practices when using tools like MT and ChatGPT in their
jobs.
\end{itemize}

Platform correctness---whether these systems are beneficial, safe, fair, and
aligned with human purposes---ultimately resides in this tacit and
institutional space, not in any single metric inside (E_{\mathrm{platform}}).

\subsection{Large-Scale Recommender Systems}

\subsubsection*{Z.12.1 Overview}

Modern recommender systems drive:

\begin{itemize}
\item video and content discovery (YouTube, TikTok, Netflix);
\item product discovery and advertising (Amazon, e-commerce, app stores);
\item social feeds and news ranking (various platforms).
\end{itemize}

They have evolved from nearest-neighbour and matrix factorisation methods to
deep, sequence-based and graph-based models that learn rich user/item
representations.

Explicit objectives typically blend:

\begin{itemize}
\item short-term engagement (clicks, watch time);
\item long-term retention and revenue;
\item sometimes additional constraints (diversity, safety rules,
content policies).
\end{itemize}

Field experiments show that recommender changes can significantly affect
what people consume, with ambiguous effects on diversity and welfare.

\subsubsection*{Z.12.2 FRFP view: explicit pipelines and population-level tacit states}

From an FRFP standpoint:

\begin{itemize}
\item The explicit mapping
(
f_{\mathrm{rec}} : \text{context, history} \to \text{ranked items}
)
is optimised for platform metrics, not directly for user welfare or
epistemic quality.
\item Users’ tacit states are updated as they consume recommended content:
beliefs change, habits crystallise, preferences drift; creators and
advertisers adapt their behaviour in response to what the system
rewards.
\item The platform’s institutional operators (product teams, policy
teams, regulators) shape objective functions and constraints, but
often through explicit metrics rather than direct control of tacit
outcomes.
\end{itemize}

This is a classic FRFP configuration where:

\begin{itemize}
\item explicit systems operate at massive scale and high frequency;
\item tacit and institutional effects (on public discourse, attention,
markets) accumulate over long horizons;
\item governance is largely metric-driven, risking explicit-only control
over complex population-level phenomena.
\end{itemize}

\subsubsection*{Z.12.3 Boundaries and safe horizons}

Key boundary questions:

\begin{itemize}
\item How much of a user’s informational diet should the recommender
effectively control?
\item To what extent are editorial and civic responsibilities delegated to
optimisation over engagement metrics?
\item How are long-horizon externalities (polarisation, addiction, market
power) fed back into objective design?
\end{itemize}

FRFP would recommend:

\begin{itemize}
\item explicit guardrails and quotas (e.g.\ diversity constraints,
exploration budgets, limits on reinforcement of recent behaviour);
\item safe horizons in user modelling (e.g.\ bounded memory, reset
mechanisms, opportunities for users to re-declare goals and
interests);
\item institutional oversight bodies that treat recommender optimisation as
a socio-technical process, not purely a metric-maximisation problem.
\end{itemize}

Current practice partially reflects this (e.g.\ “responsible recommender”
research on diversity and sustainability), but FRFP suggests a more formal
architectural treatment.

\subsection{Neural Machine Translation (NMT)}

\subsubsection*{Z.12.4 Overview}

Google Neural Machine Translation (GNMT) and successors replaced phrase-
based MT with deep neural sequence-to-sequence models using attention, and
later Transformer-based models, enabling more fluent, context-aware
translations.

By the late 2010s:

\begin{itemize}
\item Google Translate supported 100+ languages, many via a single
multi-lingual NMT model;
\item NMT reduced error rates substantially (e.g.\ claims of $\sim$60%
error reduction on some benchmarks when first deployed);
\item quality has continued to improve, particularly for high-resource
language pairs.
\end{itemize}

NMT is now woven into browsers, phones, and productivity apps, enabling
cross-lingual communication at global scale.

\subsubsection*{Z.12.5 FRFP view: explicit semantics with clear but limited ground truth}

NMT is closer to AlphaFold than to recommenders in one sense:

\begin{itemize}
\item the task---map a source sentence to an adequate target translation---
has relatively clear explicit semantics, even though nuance and
pragmatics complicate evaluation;
\item there is abundant training data (parallel corpora) and usable
automatic metrics (BLEU, COMET, etc.) as proxies for human quality;
\item human evaluation remains the ultimate tacit arbiter, especially in
professional or high-stakes translation.
\end{itemize}

FRFP interpretation:

\begin{itemize}
\item the mapping
(
f_{\mathrm{NMT}} : \text{source text} \to \text{target text}
)
has strong but not perfect explicit semantics; some meaning,
politeness, and cultural nuance lives irreducibly in tacit space;
\item everyday use (tourism, casual chat, rough reading) often places
correctness de facto in the model, with users accepting outputs
without deep scrutiny;
\item professional translators treat NMT as a drafting tool, keeping final
correctness in tacit judgement---an FRFP-aligned boundary.
\end{itemize}

\subsubsection*{Z.12.6 Boundaries by use case}

FRFP would distinguish:

\begin{itemize}
\item \textbf{Low-stakes, everyday use}:
using NMT to get the gist of a foreign-language web page; explicit
semantics are good enough; tacit oversight is light but acceptable.
\item \textbf{High-stakes or formal contexts}:
legal contracts, medical instructions, diplomatic communication;
here, NMT outputs should be treated as drafts needing human
verification, not final artefacts.
\end{itemize}

Misuse arises when the boundary is blurred---for example, when institutions
rely on raw MT output for official communications without adequate review.
FRFP suggests organisational policies that encode these boundaries: which
uses can be automated, which require human sign-off, and how safe horizons
are maintained (e.g.\ when domain drift or new jargon makes previous
evaluation obsolete).

\subsection{ChatGPT and General-Purpose LLMs}

\subsubsection*{Z.12.7 Scale and role}

ChatGPT-style LLMs are now general-purpose language tools:

\begin{itemize}
\item OpenAI reports $\sim$700--800~million weekly active users in 2025
and 2–2.5~billion prompts per day.
\item Usage spans everyday personal tasks, education, writing, coding, and
work support; a recent study finds $\sim$75% of consumer usage is
practical guidance, information seeking, and writing.
\end{itemize}

Unlike recommenders or MT, LLMs are:

\begin{itemize}
\item open-domain;
\item generative rather than purely translational or ranking-based;
\item increasingly integrated into tools that can act (agents, office
suites, browsers).
\end{itemize}

\subsubsection*{Z.12.8 FRFP view: extremely broad explicit semantics, strong tacit coupling}

From an FRFP standpoint:

\begin{itemize}
\item the explicit mapping
(
f_{\mathrm{LLM}} : \text{prompt, context} \to \text{text (and actions)}
)
is learned from huge corpora and reflects statistical structure in
language use;
\item unlike NMT or DR screening, there is no single task or metric:
correctness is highly context-dependent and often normative;
\item hallucination, bias, and value misalignment are structural, not
incidental, features of this explicit mapping.
\end{itemize}

This makes the explicit/tacit split especially delicate:

\begin{itemize}
\item users can easily over-assign correctness authority to the system
because it is fluent and broadly knowledgeable;
\item tacit skills (search, critical reading, drafting) may atrophy if
large portions of cognitive labour are repeatedly outsourced;
\item organisational processes may start to treat LLM outputs as de facto
truth or policy guidance, even where tacit judgement is essential.
\end{itemize}

\subsubsection*{Z.12.9 Boundaries and safe horizons for LLM use}

FRFP suggests several structural design choices:

\begin{itemize}
\item \textbf{Role labelling}:
treat LLMs as assistants (drafting, summarisation, brainstorming),
not as oracles or autonomous decision-makers, except in very narrow,
well-specified domains.
\item \textbf{Human-in-the-loop boundaries}:
require human endorsement for filings, public statements, and
high-stakes decisions drafted with LLM assistance (cf.\ the \emph{Mata
v.\ Avianca} case study).
\item \textbf{Safe horizons in workflows}:
limit how far LLMs can drive multi-step processes without human
checkpoints (e.g.\ code deployment, legal filings, medical advice);
design agent architectures where each major boundary event is
explicitly approved.
\item \textbf{Institutional oversight}:
treat adoption of LLMs in organisations as an FRFP-level change,
requiring policies, training, and governance, not just tool rollout.
\end{itemize}

At the population level, FRFP also worries about:

\begin{itemize}
\item epistemic effects (LLMs as a default interface to knowledge);
\item labour and skill effects (who gains or loses tacit capabilities);
\item cultural and linguistic homogenisation (if models favour high-
resource language norms).
\end{itemize}

These require institutional operators (educational systems, media, policy)
to treat LLM deployment as a long-horizon socio-technical transformation,
not just a productivity boost.

\subsection{Cross-Cutting FRFP Lessons from Platform AI}

Taken together, recommender systems, NMT, and ChatGPT / LLMs illustrate:

\begin{enumerate}
\item \textbf{Explicit strength is uneven across tasks.}
NMT has relatively clear semantics and ground truth; recommenders
and LLMs operate in domains where correctness is entangled with
tacit and normative judgements.

\item \textbf{Population-level effects dominate.}
These systems act not just on individual users but on entire
populations’ attention, beliefs, and workflows; FRFP’s population-
level semantics and institutional operators are essential to reason
about them.

\item \textbf{Boundaries are often informal or implicit.}
Everyday users rarely have clear cues about when to trust or distrust
MT and LLM outputs, or about the objectives behind recommenders.
FRFP argues for explicit boundary design and user-facing semantics.

\item \textbf{Safe horizons are rarely explicit but deeply needed.}
Long-horizon reliance on recommenders and LLMs without reset or
re-grounding mechanisms risks drift in tacit skills, norms, and
epistemics.
NMT is somewhat safer here because tasks are local and evaluation is
easier, but high-stakes use still needs boundaries.

\item \textbf{Institutional non-automation is central.}
Platform governance, media regulation, and organisational policy
cannot be replaced by optimisation over engagement or proxy metrics;
they require FRFP-style tacit judgement and multi-stakeholder
oversight.
\end{enumerate}

These platform-scale successes thus function in the FRFP monograph as
counterpoints to purely “failure” case studies:
they show that explicit AI can be enormously powerful and useful, but also
how much of its real-world impact is governed---for better or worse---by the
tacit and institutional semantics in which it is embedded.
% ============================================================
% Retrospective FRFP Case Study (Public Health / Industry):
% Google Flu Trends — A Run-Level Failure of “Big Data” Surveillance
% (written in the same two-level audit discipline used for London Whale)
% ============================================================

\section{Retrospective FRFP Case Study: Google Flu Trends (GFT)}
\label{sec:frfp-gft}

Google Flu Trends (GFT) was an influential attempt to use search-query data to estimate influenza-like illness (ILI) activity in near real time. It launched with promising early results and became a widely cited example of “big data” outperforming traditional surveillance \cite{Ginsberg2009Nature}. Its later, highly visible overestimates—especially during the 2012–2013 season—made it a canonical case where a system that looked strong in snapshots failed as a long-horizon, institution-embedded workflow \cite{Butler2013Nature,Lazer2014Science}. The FRFP lens treats this not as a single model error but as a run-level failure: feedback loops, boundary capture, weak re-grounding, and governance that drifted toward explicit-only proxies.

\subsection{Background and scope (what is the unit of analysis?)}

\paragraph{Case unit.}
We study GFT as a socio-technical program: a public-facing surveillance signal produced by a proprietary pipeline and consumed by external institutions (media, researchers, and implicitly the public health ecosystem). The unit is not “a model,” but the \emph{end-to-end workflow} from search behavior to published estimates and downstream interpretation.

\paragraph{Reference baseline.}
Traditional influenza surveillance in the United States includes ILINet, which monitors outpatient visits for ILI/respiratory illness and reports them through CDC FluView \cite{CDCILINetOverview}.

\paragraph{Key public record anchors.}
GFT’s original approach is documented in the early publication describing query-based ILI estimation \cite{Ginsberg2009Nature}.
The most widely cited retrospective critique of GFT’s failure dynamics is “The Parable of Google Flu” \cite{Lazer2014Science}.
Google later discontinued public GFT reporting \cite{Google2015NextChapterFluTrends}.

\subsection{Two-level audit discipline}

We apply a two-level audit: (L1) the \emph{system-and-operator layer} (the pipeline, its updates, and its internal controls) and (L2) the \emph{institutional-and-boundary layer} (how the signal was endorsed, amplified, constrained, and corrected).

\subsubsection*{L1: System-and-operator audit (pipeline, updates, internal controls)}

\paragraph{Explicit artefact ledger (what the system manipulates).}
Key explicit artefacts included:
\begin{itemize}
  \item query logs and derived query features;
  \item the model (or mapping) from query patterns to an ILI estimate;
  \item published time series outputs (weekly estimates);
  \item internal tuning choices (feature sets, weights, filtering, retraining cadence);
  \item internal evaluation reports comparing estimates to a reference series (often CDC ILI).
\end{itemize}
The crucial FRFP point: the “product” was not just a prediction but a \emph{repeatedly updated pipeline} whose outputs were treated as an external indicator \cite{Lazer2014Science}.

\paragraph{Pipeline shape and navigation (how the workflow evolves).}
GFT’s operational workflow behaved like a \emph{looping pipeline}:
\begin{itemize}
  \item repeated weekly inference (new queries $\rightarrow$ new estimate),
  \item periodic model revisions or recalibration,
  \item implicit drift as search behavior changes (language, media effects, interface changes),
  \item continued publication that increases reliance and visibility.
\end{itemize}
This loop structure matters because small per-period distortions can compound into large run-level divergence \cite{Lazer2014Science}.

\paragraph{Admissibility and “boundary leaks” at L1.}
The core L1 vulnerability was that the pipeline’s effective semantics depended on a stable relationship between search behavior and ILI incidence. That relationship is not static: it shifts with media attention, Google’s own product changes, and user behavior. The parable analysis emphasizes two pressure sources:
\begin{itemize}
  \item \emph{algorithmic drift}: the world changes (search behavior and platform effects),
  \item \emph{model staleness / opacity}: limited transparency about updates and validation.
\end{itemize}
These are classic long-horizon hazards: even if each week’s estimate looks plausible, the run can drift away from the intended target \cite{Lazer2014Science,Olson2013PLOSCompBio}.

\paragraph{Horizon and degradation dynamics (what made error more likely over time).}
A run-level hazard arises when:
\begin{itemize}
  \item the system is repeatedly used without a strong, enforced re-grounding cadence,
  \item the “input distribution” changes (search terms and patterns evolve),
  \item validation is partially internal and not institutionally binding.
\end{itemize}
In GFT, overestimation episodes illustrate what happens when the pipeline’s implicit assumptions degrade below the domain’s requirement floor (reliable surveillance under shifting behavior), but continued publication delays hard resets \cite{Butler2013Nature,Lazer2014Science}.

\subsubsection*{L2: Institutional-and-boundary audit (endorsement, amplification, correction)}

\paragraph{Tacit authority map (who owns “correctness”).}
No single actor had exclusive authority to declare the GFT signal “fit for use” in the way a regulated clinical tool might. Instead:
\begin{itemize}
  \item Google controlled the pipeline and publication.
  \item Public health practice relied on established surveillance (e.g., ILINet/FluView) as the authoritative reference stream \cite{CDCILINetOverview}.
  \item Media, analysts, and some researchers treated GFT as a credible near-real-time indicator, especially early on \cite{Butler2013Nature}.
\end{itemize}
This diffused authority is itself a governance risk: it encourages “soft endorsement” without explicit boundary accountability.

\paragraph{Boundary register (where the system became real).}
GFT’s boundary events were not clinical orders or trades; they were \emph{moments of institutional uptake}:
\begin{itemize}
  \item public presentation of GFT as a credible surveillance measure;
  \item third-party citation in reporting, research, and public narratives;
  \item implicit policy and attention shifts when GFT diverged from official estimates.
\end{itemize}
FRFP calls this \emph{boundary capture} risk: when an explicit output (a time series estimate) becomes treated as quasi-authoritative without a designed endorsement event that forces evidence review \cite{Lazer2014Science}.

\paragraph{Institutional mechanism sheet (how judgments were aggregated).}
The “aggregator” here was informal and media-driven:
\begin{itemize}
  \item early success + prestige of Google $\rightarrow$ increased trust,
  \item repeated citation $\rightarrow$ stabilization of belief (“this works”),
  \item later divergence $\rightarrow$ rapid swing to disillusionment.
\end{itemize}
The parable paper frames this as “big data hubris” and emphasizes that institutional narratives can lock in before the long-horizon dynamics are understood \cite{Lazer2014Science,Butler2013Nature}.

\paragraph{Auditability and replay test (could an independent reviewer reconstruct why the output changed?).}
A major L2 fragility was limited external replayability:
\begin{itemize}
  \item the full feature set and model updates were not fully public;
  \item the relevant upstream artefact (raw query stream) is not reproducible for outsiders;
  \item thus, external institutions could compare outputs to CDC data but could not reconstruct the internal causal story of changes.
\end{itemize}
This weakens governance: institutions can detect divergence, but they cannot reliably diagnose which intervention (retraining, feature changes, drift correction) would fix it \cite{Lazer2014Science}.

\paragraph{No-go pressure (what could not be guaranteed).}
GFT illustrates a general limitation: an explicit-only signal derived from behavior can fail when the behavior is endogenous to the platform and the surrounding information environment. You cannot guarantee stable semantics (search $\rightarrow$ illness) without strong, ongoing re-grounding and change control, because the measurement channel is part of the system’s own ecosystem \cite{Lazer2014Science}.

\subsection{Counterfactual: what an FRFP-aligned surveillance protocol would have required}

A counterfactual FRFP-aligned GFT would not merely “improve the model.” It would redesign boundaries, horizons, and governance:

\begin{itemize}
  \item \textbf{Explicit boundary events for adoption.}
        Before treating GFT as a surveillance input, require an explicit endorsement step by a named public-health authority with a documented evidence bundle (backtests across seasons, subgroup error analyses, divergence response plan).

  \item \textbf{Enforced safe horizons.}
        Publish only under a time-limited certification window (e.g., quarterly), with automatic downgrade or suspension if divergence thresholds trigger (relative to ILINet/FluView and other indicators) \cite{CDCILINetOverview}.

  \item \textbf{Change-control that is institution-facing.}
        Any material pipeline change (features, weights, filtering, platform changes expected to alter search behavior) triggers a “re-endorsement required” boundary, not a silent update.

  \item \textbf{Replayable audit artefacts.}
        Provide an external audit interface: versioned model cards, change logs, pre-registered evaluation slices, and a reproducible proxy dataset sufficient to validate directional claims (even if raw queries remain private).

  \item \textbf{Explicit limits on narrative claims.}
        Require communication discipline: position outputs as a provisional indicator with known failure modes, not as replacement for established surveillance—reducing boundary capture via hype.
\end{itemize}

\subsection{What FRFP adds relative to a standard postmortem}

A non-FRFP critique often concludes: “the model drifted; big data is fragile; validate more.” FRFP adds three concrete upgrades:

\begin{enumerate}
  \item \textbf{It names the real system: a long-horizon run, not a snapshot model.}
        The looping nature of repeated publication is treated as the primary object of risk.

  \item \textbf{It forces boundary accounting.}
        The key failure is not only prediction error but \emph{how an output becomes treated as authoritative} without an explicit endorsement boundary and evidence bundle.

  \item \textbf{It makes horizons and governance enforceable.}
        “Monitor and update” becomes caps, triggers, downgrade modes, and change-control operators—things you can actually implement and audit.
\end{enumerate}

\paragraph{Outcome.}
Under FRFP, Google Flu Trends becomes a clean example of a common pattern: an explicit indicator that looks strong early, is absorbed by institutions through soft boundaries, then fails run-level stability because the measurement channel is endogenous and governance lacks enforceable horizons and replayable audit artefacts \cite{Lazer2014Science,Google2015NextChapterFluTrends}.
